\subsection{Case 2608.15963: mathematical statistics inside a security claim}\label{case:contrib-15963}

This quantum-security slice tests dependency-sensitive pruning. The SSB statistic defines an empirical support fraction, while the moment statement additionally requires that the support indicators are iid Bernoulli variables. The fixture represents the latter as a dependent theorem with an explicit premise rather than as derivable exposition. Both claims remain visible in every generated view.

The key-verification and noiseless-completeness statements cannot be reduced to probability identities. They depend on who possesses the hidden Clifford circuit, what can be computed efficiently, the presentation model, and the ideal-sampling assumption. These actor and access conditions form a semantic envelope that mathematical normalization covers only partially.

The uniform-random baseline is the $\eta=1$ specialization of the depolarizing-mixture formula, not conversely, but the claims answer different experimental questions. Both remain visible in a benchmark or security audit. This case demonstrates that dependency orientation and a view policy, not an intrinsic ``redundant'' flag, determine what can be suppressed. It also shows why assumed mathematics must not certify a cryptographic threat model.